% GENERATED FILE. Edit canonical lessons, then run npm run generate:books.
% canonical-source-hash: fnv1a64:5ad8092cf0741c5e
% canonical-lessons: RU-C316-karnaval, RU-C316-maskarad, RU-C316-parad, RU-C316-shestviye, RU-C316-yarmarka

\chapter{Carnival and parade}
\label{ch:ru-carnival-and-parade}
\hlchaptermodality{316}

\begin{chapteropening}
\textbf{By the end of this chapter:} \emph{I can talk about a carnival, a masquerade, a parade, a procession and a fair.}

Complete the last lesson of chapter 316: I can talk about a carnival, a masquerade, a parade, a procession and a fair.
\end{chapteropening}

\section[\texorpdfstring{\ru{карнавал} (karnavál)}{карнавал (karnavál)}]{\ru{карнавал} (karnavál) — a carnival}

\label{lesson:RU-C316-karnaval}



\begin{quote}
Before the new one: say the Russian for to clink glasses, then the Russian for champagne.
\end{quote}

\subsection*{You'll want to know: \ru{карнавал}}
\textbf{\ru{карнавал}} (\emph{karnavál}) — \textquotedblleft{}a carnival\textquotedblright{}.

\begin{quote}
\ru{Летом} \ru{в} \ru{городе} \ru{был} \ru{карнавал}.
\end{quote}

\begin{quote}
\emph{Létom v górode byl karnavál.}
\end{quote}

\begin{quote}
In the summer there was a carnival in the town.
\end{quote}

Masculine.

\begin{cousinweb}
\emph{\ru{карнавал}} is usually traced through French \emph{carnaval} to Italian \emph{carnevale}; EtymDB lists the Russian word with no parent.
\end{cousinweb}

\begin{grammarlens}[title={What you've built}]
One more word for telling, explaining and describing at B1.
\end{grammarlens}

\subsection*{Your turn}
\begin{itemize}
\raggedright
  \item \emph{Say it:} \emph{karnavál}
  \item \emph{Say it:} \emph{karnavál}, once more
  \item \emph{Say it:} a sentence of your own with \emph{karnavál}
  \item \emph{Recall:} say the Russian for to clink glasses, then the Russian for champagne, then say \emph{karnavál} again
\end{itemize}

\begin{tcolorbox}[breakable,colback=teal!4,colframe=teal!35!black,title={Before you move on}]
What does \ru{карнавал} mean? (\textquotedblleft{}A carnival\textquotedblright{}.) Say it once more.
\end{tcolorbox}
\section[\texorpdfstring{\ru{маскарад} (maskarád)}{маскарад (maskarád)}]{\ru{маскарад} (maskarád) — a masquerade, a fancy-dress ball}

\label{lesson:RU-C316-maskarad}



\begin{quote}
Before the new one: say the Russian for champagne, then the Russian for a carnival.
\end{quote}

\subsection*{You'll want to know: \ru{маскарад}}
\textbf{\ru{маскарад}} (\emph{maskarád}) — \textquotedblleft{}a masquerade, a fancy-dress ball\textquotedblright{}.

\begin{quote}
\ru{На} \ru{маскараде} \ru{все} \ru{были} \ru{в} \ru{масках}.
\end{quote}

\begin{quote}
\emph{Na maskaráde vse býli v máskakh.}
\end{quote}

\begin{quote}
At the masquerade everyone wore masks.
\end{quote}

Masculine.

\begin{cousinweb}
\emph{\ru{маскарад}} is usually traced through French \emph{mascarade} to Italian, from the same source as \emph{\ru{маска}} \textquotedblleft{}a mask\textquotedblright{}; EtymDB lists no parent for either.
\end{cousinweb}

\begin{grammarlens}[title={What you've built}]
One more word for telling, explaining and describing at B1.
\end{grammarlens}

\subsection*{Your turn}
\begin{itemize}
\raggedright
  \item \emph{Say it:} \emph{maskarád}
  \item \emph{Say it:} \emph{maskarád}, once more
  \item \emph{Say it:} a sentence of your own with \emph{maskarád}
  \item \emph{Recall:} say the Russian for champagne, then the Russian for a carnival, then say \emph{maskarád} again
\end{itemize}

\begin{tcolorbox}[breakable,colback=teal!4,colframe=teal!35!black,title={Before you move on}]
What does \ru{маскарад} mean? (\textquotedblleft{}A masquerade, a fancy-dress ball\textquotedblright{}.) Say it once more.
\end{tcolorbox}
\section[\texorpdfstring{\ru{парад} (parád)}{парад (parád)}]{\ru{парад} (parád) — a parade}

\label{lesson:RU-C316-parad}



\begin{quote}
Before the new one: say the Russian for a carnival, then the Russian for a masquerade.
\end{quote}

\subsection*{You'll want to know: \ru{парад}}
\textbf{\ru{парад}} (\emph{parád}) — \textquotedblleft{}a parade\textquotedblright{}.

\begin{quote}
\ru{Мы} \ru{смотрели} \ru{парад} \ru{на} \ru{главной} \ru{улице}.
\end{quote}

\begin{quote}
\emph{My smotréli parád na glávnoy úlitse.}
\end{quote}

\begin{quote}
We watched the parade on the main street.
\end{quote}

Masculine.

\begin{cousinweb}
\emph{\ru{парад}} is usually traced through French or German \emph{Parade} to a Romance word built on Latin \emph{parāre} \textquotedblleft{}to prepare\textquotedblright{} (Lewis \& Short); EtymDB lists no parent.
\end{cousinweb}

\begin{grammarlens}[title={What you've built}]
One more word for telling, explaining and describing at B1.
\end{grammarlens}

\subsection*{Your turn}
\begin{itemize}
\raggedright
  \item \emph{Say it:} \emph{parád}
  \item \emph{Say it:} \emph{parád}, once more
  \item \emph{Say it:} a sentence of your own with \emph{parád}
  \item \emph{Recall:} say the Russian for a carnival, then the Russian for a masquerade, then say \emph{parád} again
\end{itemize}

\begin{tcolorbox}[breakable,colback=teal!4,colframe=teal!35!black,title={Before you move on}]
What does \ru{парад} mean? (\textquotedblleft{}A parade\textquotedblright{}.) Say it once more.
\end{tcolorbox}
\section[\texorpdfstring{\ru{шествие} (shéstviye)}{шествие (shéstviye)}]{\ru{шествие} (shéstviye) — a procession}

\label{lesson:RU-C316-shestviye}



\begin{quote}
Before the new one: say the Russian for a masquerade, then the Russian for a parade.
\end{quote}

\subsection*{You'll want to know: \ru{шествие}}
\textbf{\ru{шествие}} (\emph{shéstviye}) — \textquotedblleft{}a procession\textquotedblright{}.

\begin{quote}
\ru{Шествие} \ru{шло} \ru{по} \ru{центру} \ru{города} \ru{два} \ru{часа}.
\end{quote}

\begin{quote}
\emph{Shéstviye shlo po tséntru góroda dva chasá.}
\end{quote}

\begin{quote}
The procession moved through the town centre for two hours.
\end{quote}

Neuter.

\begin{cousinweb}
\emph{\ru{шествие}} is a Church Slavonic word, \emph{šĭstvie} \textquotedblleft{}a going, a procession\textquotedblright{} (EtymDB), built on \emph{\ru{шествовать}} \textquotedblleft{}to walk solemnly\textquotedblright{}.
\end{cousinweb}

\begin{grammarlens}[title={What you've built}]
One more word for telling, explaining and describing at B1.
\end{grammarlens}

\subsection*{Your turn}
\begin{itemize}
\raggedright
  \item \emph{Say it:} \emph{shéstviye}
  \item \emph{Say it:} \emph{shéstviye}, once more
  \item \emph{Say it:} a sentence of your own with \emph{shéstviye}
  \item \emph{Recall:} say the Russian for a masquerade, then the Russian for a parade, then say \emph{shéstviye} again
\end{itemize}

\begin{tcolorbox}[breakable,colback=teal!4,colframe=teal!35!black,title={Before you move on}]
What does \ru{шествие} mean? (\textquotedblleft{}A procession\textquotedblright{}.) Say it once more.
\end{tcolorbox}
\section[\texorpdfstring{\ru{ярмарка} (yármarka)}{ярмарка (yármarka)}]{\ru{ярмарка} (yármarka) — a fair}

\label{lesson:RU-C316-yarmarka}



\begin{quote}
Before the new one: say the Russian for a parade, then the Russian for a procession.
\end{quote}

\subsection*{You'll want to know: \ru{ярмарка}}
\textbf{\ru{ярмарка}} (\emph{yármarka}) — \textquotedblleft{}a fair\textquotedblright{}.

\begin{quote}
\ru{На} \ru{ярмарке} \ru{мы} \ru{купили} \ru{мёд}.
\end{quote}

\begin{quote}
\emph{Na yármarke my kupíli myod.}
\end{quote}

\begin{quote}
At the fair we bought honey.
\end{quote}

Feminine.

\begin{cousinweb}
\emph{\ru{ярмарка}} was borrowed through Polish \emph{jarmark} from German \emph{Jahrmarkt} \textquotedblleft{}a yearly market\textquotedblright{} (EtymDB), with \emph{Jahr} \textquotedblleft{}year\textquotedblright{}.
\end{cousinweb}

\begin{grammarlens}[title={What you've built}]
One more word for telling, explaining and describing at B1.
\end{grammarlens}

\subsection*{Your turn}
\begin{itemize}
\raggedright
  \item \emph{Say it:} \emph{yármarka}
  \item \emph{Say it:} \emph{yármarka}, once more
  \item \emph{Say it:} a sentence of your own with \emph{yármarka}
  \item \emph{Recall:} say the Russian for a parade, then the Russian for a procession, then say \emph{yármarka} again
\end{itemize}

\begin{tcolorbox}[breakable,colback=teal!4,colframe=teal!35!black,title={Before you move on}]
What does \ru{ярмарка} mean? (\textquotedblleft{}A fair\textquotedblright{}.) Say it once more.
\end{tcolorbox}
